本附录给出本文 2.4 节所用三维各向同性高斯核的归一化推导。
推导目标是确定常数系数 $C$ 使得核函数
\begin{equation}
G_\sigma(\mathbf{r}) = C \exp\!\left(-\frac{\|\mathbf{r}\|^2}{2\sigma^2}\right),
\qquad \mathbf{r}=(x,y,z)\in\mathbb{R}^3,
\end{equation}
满足单位质量（unit mass）条件
\begin{equation}
\int_{\mathbb{R}^3} G_\sigma(\mathbf{r})\,d\mathbf{r}=1,
\end{equation}
从而保证卷积平滑为局部加权平均且不引入整体幅值偏移。

\subsection*{A.1 一维高斯核的归一化}
首先考虑一维函数
\begin{equation}
g_\sigma(x)=c \exp\!\left(-\frac{x^2}{2\sigma^2}\right),
\qquad x\in\mathbb{R}.
\end{equation}
令其满足 $\int_{\mathbb{R}} g_\sigma(x)\,dx=1$，则
\begin{equation}
1=\int_{-\infty}^{\infty} c \exp\!\left(-\frac{x^2}{2\sigma^2}\right)\,dx
=c\int_{-\infty}^{\infty}\exp\!\left(-\frac{x^2}{2\sigma^2}\right)\,dx.
\end{equation}
作变量代换 $x=\sigma t$（$\sigma>0$），有 $dx=\sigma dt$，因此
\begin{equation}
\int_{-\infty}^{\infty}\exp\!\left(-\frac{x^2}{2\sigma^2}\right)\,dx
=\sigma\int_{-\infty}^{\infty} e^{-t^2/2}\,dt.
\end{equation}
记
\begin{equation}
I=\int_{-\infty}^{\infty} e^{-t^2/2}\,dt,
\end{equation}
则上述积分为 $\sigma I$。为确定 $I$，考虑
\begin{equation}
\begin{aligned}
I^2&=\left(\int_{-\infty}^{\infty}e^{-t^2/2}\,dt\right)
     \left(\int_{-\infty}^{\infty}e^{-s^2/2}\,ds\right)\\
&=\iint_{\mathbb{R}^2}e^{-(t^2+s^2)/2}\,dt\,ds.
\end{aligned}
\end{equation}
将二维积分换为极坐标 $t=r\cos\theta,\, s=r\sin\theta$，并使用面积元
$dt\,ds=r\,dr\,d\theta$，得到
\begin{align}
I^2
&=\int_{0}^{2\pi}\int_{0}^{\infty}e^{-r^2/2}\,r\,dr\,d\theta
=2\pi\int_{0}^{\infty}e^{-r^2/2}\,r\,dr.
\end{align}
再令 $u=r^2/2$（$du=r\,dr$），则
\begin{equation}
I^2=2\pi\int_{0}^{\infty}e^{-u}\,du
=2\pi.
\end{equation}
因此 $I=\sqrt{2\pi}$。代回得
\begin{equation}
\int_{-\infty}^{\infty}\exp\!\left(-\frac{x^2}{2\sigma^2}\right)\,dx
=\sigma\sqrt{2\pi}.
\end{equation}
从而一维归一化常数为 $c=\frac{1}{\sigma\sqrt{2\pi}}$，即
\begin{equation}
g_\sigma(x)=\frac{1}{\sqrt{2\pi}\sigma}
\exp\!\left(-\frac{x^2}{2\sigma^2}\right).
\end{equation}

\subsection*{A.2 三维高斯核的构造与归一化}
三维各向同性高斯核可写为三方向一维核的乘积（可分离性）：
\begin{equation}
G_\sigma(x,y,z)=g_\sigma(x)\,g_\sigma(y)\,g_\sigma(z).
\end{equation}
将 A.1 的 $g_\sigma$ 代入，可得
\begin{equation}
\begin{aligned}
G_\sigma(x,y,z)
&=\left(\frac{1}{\sqrt{2\pi}\sigma}\right)^3
\exp\!\left(-\frac{x^2+y^2+z^2}{2\sigma^2}\right) \\
&=\frac{1}{(2\pi)^{3/2}\sigma^3}
\exp\!\left(-\frac{\|\mathbf{r}\|^2}{2\sigma^2}\right)\\
&=\frac{1}{(2\pi\sigma^2)^{3/2}}
\exp\!\left(-\frac{\|\mathbf{r}\|^2}{2\sigma^2}\right).
\end{aligned}
\end{equation}
下面验证其满足单位质量条件。由于 $G_\sigma$ 已写成三方向乘积，因此
\begin{equation}
\begin{aligned}
&\int_{\mathbb{R}^3}G_\sigma(x,y,z)\,dx\,dy\,dz\\
=&\left(\int_{-\infty}^{\infty}g_\sigma(x)\,dx\right)
  \left(\int_{-\infty}^{\infty}g_\sigma(y)\,dy\right)
  \left(\int_{-\infty}^{\infty}g_\sigma(z)\,dz\right) \\
=&1\cdot 1\cdot 1 = 1.
\end{aligned}
\end{equation}
因此上述系数即为三维高斯核的归一化常数。
在本文方法中，该归一化保证了平滑操作为加权平均：对常数场 $F(\mathbf{x})\equiv c$，
有 $(G_\sigma*F)(\mathbf{x})\equiv c$，从而不会引入整体幅值偏移。